\scanpage{68}

donné, section d'un faisceau structural, d'ordre infini sur chaque fibre, on trouve comme anneau d'invariants $k[\sigma_2]$ \dots.

\uline{Corollaire} Les fonctions polynômiales de degré $\leq2$ invariantes par $D_\zeta=\operatorname{Aut}(\Pi_{\zeta,k})$ sont les fonctions

\[axy+b\qquad (a,b\in k)\]

Donc à scalaire près une unique forme quadratique \unsure{invar.|invariante}, \inserted{$q(x,y)=\deleted{\unsure{q}}\sigma_2(x,y)=xy$} qu'on peut normaliser par $\sigma_2(\unsure{S|s})=1$ si $s\in S$. \uline{Une unique} \uline{cônique affine} ``circonscrite à $\Pi$'' i.e. contenant les sommets, savoir celle d'équation

\margin{\unsure{vif !|vérif !|réf !|Kif !} Ne marche que si $n\geq5$, avec démonstration à donner. Ici on a néanmoins une cônique circonscrite canonique -- invariante par $G$}
% slanted

\[xy-1=0\qquad\deleted{\unsure{xy-1=0|xy-1=2}}\]

\deleted{La forme quadratique $\sigma_2$}

Pour un $n$-polygône régulier quelconque (sur $k$ alg. clos de car première à $n$) -- pas épinglé -- il y a donc \inserted{à scalaire près} une unique forme quadratique invariante $\unsure{\zeta_0|q_0}$, qu'on normalise par

\[q_{0}(s)=1\qquad(s\in S)\]

% archive sheet number: 67
